\documentclass[11pt,a4paper]{article}
\usepackage[T1]{fontenc}
\usepackage[utf8]{inputenc}
\usepackage{lmodern}
\usepackage[french]{babel}
\usepackage[a4paper,margin=2.3cm]{geometry}
\usepackage{microtype}\usepackage{xcolor}\usepackage{enumitem}\usepackage{titlesec}\usepackage{hyperref}
\usepackage[most]{tcolorbox}

\definecolor{ink}{HTML}{1c2330}\definecolor{accent}{HTML}{2f6e4c}\definecolor{accent2}{HTML}{8a5a1c}
\definecolor{soft}{HTML}{eef3ef}\definecolor{softline}{HTML}{cfe0d4}
\definecolor{open}{HTML}{1c4e7a}\definecolor{opensoft}{HTML}{eaf1f8}
\definecolor{warn}{HTML}{7a2418}\definecolor{warnsoft}{HTML}{fbeeea}
\definecolor{data}{HTML}{5a2f6e}\definecolor{datasoft}{HTML}{f2ecf6}

\hypersetup{colorlinks=true,linkcolor=accent,urlcolor=accent2,pdftitle={Plan --- génération procédurale}}
\titleformat{\section}{\Large\bfseries\color{ink}}{\thesection.}{0.5em}{}
\titleformat{\subsection}{\large\bfseries\color{accent}}{\thesubsection}{0.5em}{}
\titlespacing{\section}{0pt}{1.2em}{0.4em}
\setlist[itemize]{leftmargin=1.2em,itemsep=2pt,topsep=2pt}
\newtcolorbox{databox}{enhanced,breakable,colback=datasoft,colframe=data,boxrule=0.7pt,left=9pt,right=9pt,top=6pt,bottom=6pt,arc=2pt,fonttitle=\bfseries,coltitle=data,title=Données XCOM 2}
\newtcolorbox{openbox}{enhanced,breakable,colback=opensoft,colframe=open,boxrule=0.5pt,left=9pt,right=9pt,top=6pt,bottom=6pt,arc=2pt}
\newtcolorbox{keybox}{enhanced,breakable,colback=soft,colframe=accent,boxrule=0.7pt,left=9pt,right=9pt,top=6pt,bottom=6pt,arc=2pt}

\title{\vspace{-1.5em}\textbf{\color{ink}\Huge Génération procédurale}\\[0.3em]
\large\color{accent}Maps \& missions --- plan d'exploration\\[0.2em]
\normalsize\color{ink}\itshape Plusieurs techniques, ancrées sur notre maillage et les données XCOM, puis auto-évaluées}
\author{}\date{}
\begin{document}\maketitle\vspace{-3em}

\section{Pourquoi (cadre DIRECTION)}
La DIRECTION choisit \textbf{A6 : procédural borné + golden path} --- variété sans perdre le contrôle, pas du «~tout aléatoire~». On a déjà la brique géométrique (\texttt{genMesh} $\to$ maillage irrégulier en île) ; il manque de \textbf{peupler} ce maillage (couvert, relief, pods, objectifs) \textbf{et} d'émettre la \textbf{mission} correspondante (slots de déploiement, ennemis, conditions, dialogues).

\section{Ce que XCOM 2 nous apprend}
\begin{databox}
\textbf{Système Plot + Parcel + PCP :}
\begin{itemize}
\item \textbf{Plot} : trame \textbf{faite main} (routes/chemins avec \emph{sockets}), commune à plein de missions $\to$ lisibilité.
\item \textbf{Parcels} : modules (petit/moyen/grand, en \textbf{tuiles}) tirés au sort et encastrés dans les sockets.
\item \textbf{PCP} (\emph{Plot Cover Parcels}) : remplissage, \textbf{$\sim$8--16 tuiles}, largeur de rue.
\item \textbf{Tuile} = \textbf{96 unités Unreal}. \textbf{Asset swapping} par biome (tempéré $\to$ aride/toundra).
\item Philosophie dir. artistique : \textbf{«~Random's not fun~»} --- mélange main + procédural pour la rejouabilité sans le chaos.
\end{itemize}
\textbf{Pods :} \textbf{2--3 ennemis} par pod (selon difficulté). \textbf{Nb de pods $\approx$ difficulté} : facile \textbf{2--4}, moyen \textbf{5--6}, difficile \textbf{7--8}. Patrouille puis \textbf{activation à la révélation} --- déjà chez nous (\texttt{pod}/\texttt{asleep}/réveil à vue).
\end{databox}
\noindent\emph{Confiance : plot/parcel \& philosophie = source dev (GDC 2018 «~Plot and Parcel~»). Tuile 96\,u, PCP, tailles/nb de pods = wikis/guides communautaires (dimensions exactes des parcels non publiées). À recalibrer chez nous (maillage $\neq$ grille).}

\section{Notre contexte}
\begin{itemize}
\item \textbf{Diffère :} maillage \textbf{Voronoï irrégulier} (cases de taille variable), pas une grille. On raisonne en \textbf{nombre de cellules}, \textbf{adjacence}, \textbf{distance en sauts}.
\item \textbf{Transpose :} densité de couvert, dimensionnement/placement des pods, plot fait-main (= notre golden path), séparation \textbf{déploiement / objectif / pods}.
\item \textbf{Déjà en place :} \texttt{genMesh}, le générateur de profil de terrain auto, murets (couvert d'arête), rocher (bloquant), relief, pods+sommeil, triggers, format mission complet, et \texttt{tools/balance.js} (\textbf{simulateur IA vs IA}).
\end{itemize}

\section{Les 5 couches à générer}
\begin{enumerate}
\item \textbf{Forme} --- maillage île (déjà : \texttt{genMesh}).
\item \textbf{Relief} --- champs d'élévation (collines, crêtes, cuvettes).
\item \textbf{Couvert \& terrain} --- rocher / accidenté / murets vers une \textbf{densité cible}.
\item \textbf{Pods \& déploiement} --- slots joueur, pods (taille, classe, sommeil), PNJ/otages.
\item \textbf{Objectif \& scénarisation} --- archétype, zones, dialogues, conditions, renforts.
\end{enumerate}

\section{Techniques à essayer}
\emph{Interface commune : \texttt{generate(seed, params) -> mission JSON}. On compare par couche.}
\begin{itemize}
\item \textbf{A. Scatter par densité (XCOM-like)} --- baseline : couvert/rocher/muret jusqu'à un \textbf{ratio cible}. Simple, ajustable.
\item \textbf{B. Stamping de prefabs / «~parcels~»} --- la leçon XCOM : bibliothèque de \textbf{chunks faits main} (gué, ruine-chokepoint, clairière, campement) estampillés sur des régions $\to$ variété \textbf{maîtrisée}.
\item \textbf{C. Bruit (Perlin/value)} --- élévation \& zones accidentées continues (prolonge le profil auto).
\item \textbf{D. Automates cellulaires} --- amas de couvert / lanes ouvertes organiques.
\item \textbf{E. Partition (BSP/régions)} --- salles + passages (cartes «~intérieures~»/fronts).
\item \textbf{F. WFC} --- cohérence d'adjacence des terrains (avancé, si A--E insuffisants).
\item \textbf{G. Placement par graphe (transversal)} --- sur l'adjacence : déploiement (un bord), objectif (loin), pods espacés le long des chemins, près du couvert, \textbf{hors de vue} au tour 1 (\texttt{los}/\texttt{hops} existent).
\end{itemize}

\section{Dimensionnement \& placement des pods}
\begin{itemize}
\item \textbf{Taille 2--3} (paramétrable selon la difficulté de la région).
\item \textbf{Nb de pods} : facile 2--4, moyen 5--6, difficile 7--8 \emph{(à revoir à la baisse à notre échelle : sans doute 1--4)}.
\item Pods \textbf{hors de vue} du déploiement au tour 1, \textbf{espacés} ($\geq$ K hops, anti multi-activation), un \textbf{gardien d'objectif}, certains \textbf{endormis} (\texttt{asleep}).
\item \textbf{Renforts} optionnels (tour N ou trigger). \textbf{Couvert} : assez pour que les lanes ouvertes soient un \emph{choix} --- départ $\sim$25\,\% cellules-couvert + murets sur $\sim$1 arête/6, à régler par simulation.
\end{itemize}

\section{Générer la MISSION, pas que la carte}
Sortie = un \textbf{mission JSON} déjà chargeable :
\begin{itemize}
\item \textbf{Slots de déploiement} (Phase A : cases joueur) sur un bord sûr.
\item \textbf{Ennemis} : \texttt{cls}/\texttt{pod}/\texttt{asleep} selon les pods ; PNJ/otages si «~sauvetage~».
\item \textbf{Archétypes} : élimination, sauvetage (otages = PNJ à protéger), extraction (atteindre une zone), défense (survivre N tours), sabotage (atteindre une cible) --- via \texttt{protect}/\texttt{loot}/conditions.
\item \textbf{Triggers} : intro/ambiance, révélation de pod, renforts, ouverture/fermeture de zone (modèles).
\item \textbf{Métadonnées geoscape} : difficulté + archétype + biome $\to$ la région \emph{demande} «~une mission sauvetage, moyenne~».
\end{itemize}

\section{Auto-évaluation --- garder les bonnes}
On a \texttt{tools/balance.js} + \texttt{tools/engine.js} (\textbf{sim IA vs IA headless}). Boucle :
\begin{enumerate}
\item \textbf{Générer} N cartes (techniques \& seeds variés).
\item \textbf{Simuler} chacune K fois (IA vs IA).
\item \textbf{Mesurer} : taux de victoire, nb de tours, pertes, \% couvert utilisé, accessibilité, multi-activations.
\item \textbf{Filtrer} : rejeter trop facile/dur/dégénéré ; garder «~tendu mais jouable~».
\item \textbf{Comparer les techniques} sur ces métriques.
\end{enumerate}
\begin{keybox}
L'angle fort : un générateur \textbf{mesuré}, pas seulement «~joli~». C'est ça, le procédural \textbf{borné} de la DIRECTION.
\end{keybox}

\section{Plan d'attaque (phases)}
\begin{itemize}
\item \textbf{P1 --- Baseline jouable} : A (densité) + G (graphe) + émission mission (élimination). Bouton «~Générer mission~» (seedé) dans l'éditeur.
\item \textbf{P2 --- Auto-évaluation} : brancher \texttt{balance.js} (générer$\to$simuler$\to$filtrer + métriques).
\item \textbf{P3 --- Prefabs/parcels} : B + bibliothèque de chunks faits main (curé, lisible).
\item \textbf{P4 --- Relief \& cohérence} : C/D (bruit, automates), éventuellement F (WFC).
\item \textbf{P5 --- Archétypes \& geoscape} : sauvetage/extraction/défense/sabotage + triggers ; la région du geoscape peuple ses missions seule.
\end{itemize}

\section{Décisions ouvertes}
\begin{openbox}
\begin{itemize}
\item \textbf{Échelle} : réduire les comptes de pods XCOM à notre taille (sans doute 1--4 pods).
\item \textbf{Curé vs varié} : équilibre prefabs (B) / scatter (A) --- la DIRECTION penche curé.
\item \textbf{Persistance} : missions générées \textbf{figées} (JSON sauvegardé, rejouable/partageable) ou \textbf{régénérées} par seed ? (figer colle au golden path \& à «~l'échec fait avancer~»).
\item \textbf{Thème} : faire \textbf{monter la densité/difficulté} à mesure que la carte du monde «~se déforme~».
\end{itemize}
\end{openbox}

\begin{center}\itshape Prochaine étape proposée : P1 --- baseline (scatter + pods + objectif) émettant une mission jouable.\end{center}

\vfill\footnotesize\textcolor{ink!70}{Sources : GDC 2018 «~Plot and Parcel: Procedural Level Design in XCOM 2~» (B. Hess) ; gamerant.com (procgen \& destructibilité) ; XCOM Wiki / Fandom (Pod, XCOM 2) ; guides Steam d'édition de niveaux XCOM 2. Dimensions exactes des parcels non publiées ; valeurs à recalibrer pour notre maillage.}
\end{document}
